\documentclass{beamer}

% Theme selection
\usetheme{Madrid}
\usecolortheme{default}

% Packages
\usepackage[utf8]{inputenc}
\usepackage[T1]{fontenc}
\usepackage{amsmath}
\usepackage{amssymb}
\usepackage{graphicx}
\usepackage{braket} % For Dirac notation

% Title information
\title[Quantum Bayesian Portfolio]{Optimasi Portofolio Menggunakan Quantum Bayesian Game Theory dengan VQE Ansatz}
\subtitle{Dari Ketidakpastian Ekonomi ke Solusi Fisika Kuantum}
\author{DreamyTA Team}
\date{\today}

\begin{document}

% --- BAGIAN 1: PENDAHULUAN & MASALAH ---

% Slide 1: Title
\begin{frame}
    \titlepage
\end{frame}

% Slide 2: Masalah Utama (Ekonomi)
\begin{frame}{Masalah: Ketidakpastian Pasar}
    \begin{block}{Model Klasik (Markowitz) vs Realita}
        \begin{itemize}
            \item \textbf{Asumsi Klasik:} Pasar bersifat \textit{deterministik}. Parameter (return $\mu$, risiko $\sigma$) dianggap statis berdasarkan data historis.
            \item \textbf{Realita:} Pasar penuh \textit{ketidakpastian (incomplete information)}.
            \item \textbf{Dampak:} Portofolio optimal pada data masa lalu seringkali \textbf{gagal total} saat kondisi pasar berubah drastis (Crash/Bearish).
        \end{itemize}
    \end{block}

    \begin{alertblock}{Problem Statement}
        Bagaimana kita membangun portofolio yang "tahan banting" (robust) terhadap perubahan kondisi pasar yang tidak diketahui sebelumnya?
    \end{alertblock}
\end{frame}

% Slide 3: Solusi Quantum Bayesian
\begin{frame}{Solusi: Quantum Bayesian Game Framework}
    Pendekatan multidisiplin yang menggabungkan tiga pilar:
    
    \begin{enumerate}
        \item \textbf{Game Theory (Bayesian):}
        Memodelkan interaksi strategi investor melawan "Nature" (Ketidakpastian Pasar) secara probabilistik.
        
        \item \textbf{Quantum Computing (EWL Scheme):}
        Memanfatkan superposisi dan entanglement untuk memperluas ruang strategi dan menghindari jebakan lokal minimum klasik.
        
        \item \textbf{VQE (Variational Quantum Eigensolver):}
        Algoritma hibrida untuk mencari konfigurasi portofolio paling optimal (Ground State) secara efisien.
    \end{enumerate}
\end{frame}


% --- BAGIAN 2: LANDASAN TEORI ---

% Slide 4: Landasan Teori - Classical Bayesian Game
\begin{frame}{Landasan Teori: Konsep Bayesian Game}
    Dalam Bayesian Game, pemain bertindak tanpa mengetahui "tipe" lawan secara pasti.
    
    \begin{itemize}
        \item \textbf{Players:} Investor vs Nature (Alam/Pasar).
        \item \textbf{Types of Nature ($\Theta$):}
        Alam memiliki dua kemungkinan identitas (State):
        \begin{itemize}
            \item $\theta_1$: \textbf{Bullish} (Pasar Naik)
            \item $\theta_2$: \textbf{Bearish} (Pasar Turun/Crash)
        \end{itemize}
        \item \textbf{Prior Belief ($P$):}
        Keyakinan investor terhadap kondisi pasar:
        $$ P(\text{Bull}) = q, \quad P(\text{Bear}) = 1-q $$
    \end{itemize}
\end{frame}

% Slide 5: Landasan Teori - Quantum EWL Scheme
\begin{frame}{Landasan Teori: Skema Kuantum EWL}
    EWL (Eisert-Wilkens-Lewenstein) adalah protokol untuk memainkankan game di sirkuit kuantum.
    
    \begin{center}
    \textbf{Komponen Utama:}
    \end{center}
    
    \begin{enumerate}
        \item \textbf{State Awal ($\ket{\psi_0}$):} Superposisi strategi.
        \item \textbf{Entanglement Operator ($\hat{J}$):} 
        Menciptakan korelasi kuantum antar aset yang tidak mungkin ada di fisika klasik.
        $$ \hat{J} = \exp(i \frac{\gamma}{2} \hat{\sigma}_x^{\otimes n}) $$
        \item \textbf{Strategi Operator ($\hat{U}$):}
        Pemain memutar qubit mereka untuk memilih strategi (Beli/Jual).
    \end{enumerate}
\end{frame}


% --- BAGIAN 3: METODOLOGI (CRITICAL) ---

% Slide 6: Metodologi - Langkah 1 (Matriks Payoff)
\begin{frame}{Metodologi: Konstruksi Matriks Payoff Kuantum}
    Alih-alih satu matriks, kita mendefinisikan operator Payoff (Hamiltonian) terpisah untuk setiap kondisi pasar.
    
    \begin{block}{1. Matriks Payoff Kondisi Bullish ($\hat{H}_{Bull}$)}
    $$ \hat{H}_{Bull} = \sum_{i,j} w \sigma_{ij}^{bull} \hat{Z}_i \hat{Z}_j - (1-w) \mu_i^{bull} \hat{Z}_i $$
    \footnotesize{Menggunakan kovarians dan return saat pasar naik.}
    \end{block}
    
    \begin{block}{2. Matriks Payoff Kondisi Bearish ($\hat{H}_{Bear}$)}
    $$ \hat{H}_{Bear} = \sum_{i,j} w \sigma_{ij}^{bear} \hat{Z}_i \hat{Z}_j - (1-w) \mu_i^{bear} \hat{Z}_i + \text{Penalty}_{crash} $$
    \footnotesize{Menggunakan parameter saat krisis, ditambah penalti risiko ekstrem.}
    \end{block}
\end{frame}

% Slide 7: Metodologi - Langkah 2 (Bayesian Weighted)
\begin{frame}{Metodologi: Bayesian Weighted Hamiltonian}
    Inilah inti dari pendekatan \textbf{Quantum Bayesian}. Kita menggabungkan kedua matriks payoff berdasarkan probabilitas (Belief) investor.
    
    $$ \hat{H}_{Total} = \mathbb{E}[\hat{H}] = \sum_{\theta \in \Theta} P(\theta) \cdot \hat{H}_{\theta} $$
    
    \begin{alertblock}{Derivasi Hamiltonian Akhir}
    \Large
    $$ \hat{H}_{QBG} = q \cdot \hat{H}_{Bull} + (1-q) \cdot \hat{H}_{Bear} $$
    \end{alertblock}
    
    \vspace{0.2cm}
    \textbf{Artinya:} Kita mencari portofolio yang meminimalkan energi rata-rata di seluruh kemungkinan semesta (Bull & Bear) secara simultan.
\end{frame}

% Slide 8: Metodologi - Langkah 3 (Penalti & Kendala)
\begin{frame}{Metodologi: Penambahan Kendala (Penalty Term)}
    Hamiltonian $\hat{H}_{QBG}$ belum lengkap tanpa kendala anggaran (jumlah aset = $K$).
    
    \textbf{Fungsi Penalti Lagrangian:}
    $$ \hat{H}_{constraint} = \lambda \left( \sum_{i=1}^N \frac{1-\hat{Z}_i}{2} - K \right)^2 $$
    
    \textbf{Hamiltonian Lengkap untuk VQE:}
    $$ \hat{H}_{VQE} = \hat{H}_{QBG} + \hat{H}_{constraint} $$
    
    Inilah operator hermit, berupa matriks raksasa ($2^N \times 2^N$), yang akan dicari Ground State-nya.
\end{frame}


% --- BAGIAN 4: EKSEKUSI KUANTUM ---

% Slide 9: Eksekusi - Mapping ke Qubit (Ising)
\begin{frame}{Eksekusi: Mapping ke Model Ising}
    Agar dapat dijalankan di Komputer Kuantum, variabel biner portofolio $x_i \in \{0,1\}$ dipetakan ke spin operator $\hat{Z}_i$.
    
    \begin{columns}
        \column{0.5\textwidth}
        \textbf{Domain Klasik ($x_i$):}
        \begin{itemize}
            \item $x_i = 1$: Aset Dipilih
            \item $x_i = 0$: Aset Tidak Dipilih
        \end{itemize}
        
        \column{0.5\textwidth}
        \textbf{Domain Kuantum ($\hat{Z}_i$):}
        \begin{itemize}
            \item $\Ket{1}$ (Spin Down): Terpilih
            \item $\Ket{0}$ (Spin Up): Tidak Terpilih
        \end{itemize}
    \end{columns}
    
    \vspace{0.5cm}
    Transformasi: $x_i \rightarrow \frac{1-\hat{Z}_i}{2}$
\end{frame}

% Slide 10: Eksekusi - Algoritma VQE
\begin{frame}{Eksekusi: Algoritma VQE}
    \textbf{Variational Quantum Eigensolver (VQE)} bekerja sebagai loop hibrida:
    
    \begin{enumerate}
        \item \textbf{QPU (Quantum):} Menyiapkan state ansatz $\ket{\psi(\vec{\theta})}$ dan mengukur ekspektasi energi $\braket{\psi | \hat{H}_{VQE} | \psi}$.
        \item \textbf{CPU (Klasik):} Menggunakan optimizer (COBYLA/SPSA) untuk memperbarui parameter $\vec{\theta}$ agar energi semakin kecil.
        \item \textbf{Konvergensi:} Loop berulang hingga energi minimum dicapai (Ground State).
    \end{enumerate}
\end{frame}

% Slide 11: Eksekusi - Hasil Optimasi
\begin{frame}{Eksekusi: Output Portofolio}
    Hasil akhir dari VQE adalah state kuantum dengan probabilitas tertinggi untuk diukur.
    
    \textbf{Interpretasi Hasil:}
    Jika hasil pengukuran state adalah $\ket{1101...}$, maka:
    \begin{itemize}
        \item Aset 1: Beli
        \item Aset 2: Beli
        \item Aset 3: Jangan Beli
        \item Aset 4: Beli
        \item dst.
    \end{itemize}
    
    Portofolio ini secara matematis adalah \textit{Best Response} terhadap strategi campuran Nature (Bull/Bear).
\end{frame}


% --- BAGIAN 5: KESIMPULAN ---

% Slide 12: Kesimpulan
\begin{frame}{Kesimpulan}
    \begin{enumerate}
        \item \textbf{Novelty:} Penggabungan \textbf{Bayesian Game Theory} dengan \textbf{Quantum VQE} memberikan kerangka kerja baru untuk optimasi portofolio di bawah ketidakpastian.
        
        \item \textbf{Robustness:} Dengan memisahkan matriks payoff ($\hat{H}_{Bull}$ dan $\hat{H}_{Bear}$) dan menggabungkannya via pembobotan Bayesian, solusi yang dihasilkan lebih tahan terhadap crash pasar.
        
        \item \textbf{Keunggulan Kuantum:} Skema EWL dan VQE memungkinkan eksplorasi ruang solusi yang lebih luas dibanding algoritma optimasi klasik, terutama untuk masalah kombinatorial yang kompleks.
    \end{enumerate}
\end{frame}

\end{document}
